\documentclass[10pt,a4paper]{article}
\usepackage{amsmath,amssymb,amsthm,graphicx}
\usepackage{amssymb,latexsym}
\usepackage{amsmath}
\usepackage{amsfonts}
\usepackage{amsthm}
\usepackage[colorlinks,pagebackref]{hyperref}
\newtheorem{theorem}{Theorem}%[section]
\newtheorem{problem}[theorem]{Problem}%[section]
\newtheorem{observation}[theorem]{Observation}
\newtheorem{lemma}[theorem]{Lemma}
\newtheorem{corollary}[theorem]{Corollary}
\newtheorem{definition}{Definition}
\newtheorem{example}[theorem]{Example}
\newtheorem{conjecture}{Conjecture}%[section]
\newtheorem{proposition}[theorem]{Proposition}
\newtheorem{claim}[theorem]{Claim}
\newtheorem{remark}[theorem]{Remark}
\def\qed{\hfill $\Box$}

\def\R{{\mbox{\boldmath$\mathcal R$}}}
%\def\R{{\bf\mathcal R}}
\def\F{{\mathcal F}}
\def\P{{\mathcal P}}
\title{A72-vertex perfect graph requiring three comparability subgraphs}
\author{András Gyárfás, Márton Marits, Géza Tóth}
\date{arxiv:2408.13523v2; CC BY 4.0}
\ifdefined\backrefparscanfalse\backrefparscanfalse\fi
\begin{document}
\maketitle
\noindent\textit{Editorial note.} H is the Cartesian product K9 square K4 with one new pendant vertex attached to each of its36 grid vertices; it has72 vertices. This is a finite graph, not an asymptotic construction.\par\medskip
\noindent\textit{Editorial note.} The lower bound allows A and B to overlap. Their transitive orientations are considered separately, exactly as in source667; no partition assumption is inserted into the cover proof.\par\medskip
\noindent\textit{Editorial note.} A row vertex isolated in one of the two oriented row subgraphs is both maximal and minimal. For the four-type counting argument it may be assigned either applicable status, with one fixed choice per orientation; same-type vertices still cannot share any covering edge.\par\medskip
\noindent\textit{Editorial note.} The X\_i/Y\_i mentions in source661 and X\_i in680 refer to the row and column families defined648-649; their names are dummy indices, and the author notation is retained.\par\medskip
\noindent\textit{Editorial note.} This task uses the fixed source version2408.13523v2. The official abstract page lists a later v3; no unreviewed replacement of the selected source is made.\par\medskip
\noindent\textit{Editorial scope.} One explicit finite graph and the exact values c(H)=p(H)=3, together with its perfectness. Interval graph asymptotic bounds, the earlier recursive upper bound, and the50-vertex personal-communication improvement are not counted.\par\medskip
\noindent\textit{Allowed standard prerequisites (editorial).}\begin{itemize}
\item Transitive orientations characterize comparability graphs. An orientation is transitive when u->v and v->w imply that uw is an edge and is oriented u->w. Comparability graphs are exactly graphs admitting such an orientation.. Each covering comparability graph carries its own transitive orientation; an edge belonging to both is allowed two orientations.
\item Bipartite line graphs are perfect. Every induced subgraph of the line graph of a bipartite graph has chromatic number equal to clique number, by the classical bipartite edge-coloring theorem. Attaching pendant vertices preserves perfectness.. Certify only perfectness of the explicit grid-with-leaves graph. No lower bound on the cover parameter is assumed.
\item Elementary finite pigeonhole principle. Convert four types of bad row vertices into three good vertices in one column after the nontrivial local orientation constraints are established.
\end{itemize}\medskip
Perfect graphs are the graphs $G$ in which all induced subgraphs $H\subseteq G$
satisfy the property $\omega(H)=\chi(H)$ where $\omega(H)$ and $\chi(H)$ denote
the size of the largest complete subgraph and the chromatic number.

Here we focus on a well known family of perfect graphs, the {\em comparability graphs},
which can be obtained from a partially ordered set as follows.
Vertices correspond to the elements of the partially ordered set,
and two vertices are connected by an undirected edge if and only if the corresponding elements
are comparable. The perfectness of comparability graphs are established in \cite{D50}, \cite{M71}.

Let $p(G)$ denote the minimum number $m$ such that $E(G)=\cup_{i=1}^m E(G_i)$ where the $G_i$-s are edge-disjoint comparability subgraphs of $G$.
Dropping  the condition that the $G_i$-s are edge-disjoint, we define $c(G)$ as the smallest number of comparability subgraphs of $G$ needed to {\em cover} $E(G)$.
Clearly, $p(G)\ge c(G)$ but seems not easy to find examples where the gap between them is large. For example, we do not know any perfect $G$ for which $p(G)>c(G)$.

The {\em line graph} of a graph $G$ is obtained by representing the edges of $G$ by vertices and
considering two such vertices adjacent if the corresponding edges of $G$ have a common vertex.
The line graphs of bipartite graphs, LBIP, form a well-known class of perfect graphs.

For the $c(G)=3$ case we can construct a smaller
example  based on the cartesian product  $K_n\square K_m$, the line-graph of the complete bipartite graph $K_{n,m}$.
Let $H_{9,4}$ be the graph with $72$ vertices,
obtained from $K_9\square K_4$ by attaching a pendant edge at each vertex.

\begin{proposition}\label{H9} The graph $H=H_{9,4}$ is perfect and $p(H)=c(H)=3$.
\end{proposition}

\noindent{\bf Proof.}



Let $V(H)=K\cup W$ where
$$K=\{v_{i,j}: 1\le i\le 9, 1\le j\le 4\}, W=\{w_{i,j}:  1\le i\le 9, 1\le j\le 4\}.$$
For $1\le j\le 4$, let $X_j=\{v_{i, j} : 1\le i \le 9\}$ be the {\em $j$-th row} of $K$ and
for  $1\le i\le 9$, let $Y_i=\{v_{i, j} : 1\le j \le 4\}$ be the {\em $i$-th column} of $K$.

There are three types of edges,
(i)  for  $1\le i\le 9$, $1\le j\le 4$,  $e_{i,j}=\{v_{i,j},w_{i,j}\}$ are the pendant edges of $H$,
(ii) for $1\le j\le 4$, $X_j$ induces a complete graph
(iii)  for $1\le i\le 9$, $Y_i$ also induces a complete graph.

  The graph $H$ is perfect since $K_9\square K_4$ is the line graph of the complete bipartite graph
  $K_{9,4}$ and adding pendant edges to a perfect graph preserves perfectness.

  To see that $c(H)\le p(H)\le 3$,
  note that the transitive orientations on the complete graphs spanned by the
  $X_i$-s (resp. $Y_i$-s), and any orientation on the pendant edges provide a partition of $E(H)$ into three comparability subgraphs.

  On the other hand, suppose for a contradiction that
  $H$  is the union of two comparability subgraphs, $A$ and $B$.
%defined by $\prec_A$ and $\prec_B$.
Consider $A$ and $B$
with their transitive orientations (which might introduce multiply directed edges).

A vertex $v_{i,j}\in X_j$ is  $A$-good (resp. $B$-good)
if it is
the midpoint of an oriented two-edge path in
row $X_j$ with both edges in $A$ (resp. $B$).
A vertex $v_{i,j}$ is good if it is $A$-good or $B$-good and bad otherwise.

Observe that if $v_{i,j}$ is bad, than it has to be maximal or minimal in both
transitively oriented subgraphs $A$ and $B$, restricted to $X_j$.
So there are four types of bad vertices, ($A$-max, $B$-max),
($A$-max, $B$-min),  ($A$-min, $B$-max),  ($A$-min, $B$-min).

Let $u$ and $v$ be two bad vertices
in a row $X_i$. Suppose that $uv$ belongs to $A$.
Then $u$ and $v$ cannot be $A$-max (resp. $A$-min) at the same time.
Similarly, if $uv$ belongs to $B$,
then $u$ and $v$ cannot be $B$-max (resp. $B$-min) at the same time.
Therefore, in each row $X_i$ there is at most one bad vertex of each type, so there are at most $4$ bad vertices.
Consequently, in each row there are at least $5$ good vertices.
So $K$ contains at least $20$ good vertices. It follows that one of the columns contains at least $3$ good vertices.
Assume without loss of generality that it is $Y_1$ and the three good vertices are $v_{1,1}$, $v_{1,2}$, $v_{1,3}$.

Suppose that $v_{1,1}$ is $A$-good and $v_{1,2}$ is $B$-good.
The edge $v_{1,1}v_{1,2}$ belongs to $A$ or to $B$, and in both cases, by transitivity,
we must have an edge of $H$ which is neither vertical, nor horizontal, which is a contradiction.
Therefore,  $v_{1,1}$ and $v_{1,2}$ are of the same type. Analogously, $v_{1,3}$ is also of the same type, say, all of them are $A$-good.
Consider now the edges between  $v_{1,1}$, $v_{1,2}$ and  $v_{1,3}$. If any of them belongs to $A$, then again  by transitivity,
we must have an edge of $H$ which is neither vertical, nor horizontal, which is a contradiction.
So, all three edges belong to $B$. Since $B$ has a transitive orientation,
$v_{1,1}$, $v_{1,2}$ and  $v_{1,3}$ form an oriented path in $B$. Suppose that $v_{1,1}$ is its middle vertex.
Then $v_{1,1}$ is the middle vertex of an $A$-path and $B$-path in $K$.
Finally, no matter whether $e_{1,1}$ belongs to $A$ or to $B$ and no matter which direction is assigned to it,
by the transitivity of $A$ or $B$ we must have another edge from $w_{1,1}$ which is a contradiction.
\qed.

\begin{thebibliography}{99}
\bibitem{D50} R. P.  Dilworth, A Decomposition Theorem for Partially Ordered Sets, {\em Annals of Mathematics} {\bf 51} (1950), 161--166.
\bibitem{M71} L. Mirsky, A dual of Dilworth's decomposition theorem, {\em American Mathematical Monthly} {\bf 78} (1971), 876--877.
\end{thebibliography}
\end{document}
